\section{背景與動機}
\label{sec:background}

\subsection{生產系統中的慢故障}

灰色故障（gray failure）~\cite{huang2017gray}泛指各種部分故障，其中元件並未完全停機，卻也無法正常運作。
慢故障是灰色故障的一個子類，其特徵為服務時間延長至 $s$ 倍（$s > 1$）。
常見成因包括：記憶體壓力導致過度垃圾回收、
磁碟退化增加 I/O 延遲、溫度限流降低 CPU 頻率，
或特定路徑上的網路擁塞~\cite{gunawi2018fail}。

與崩潰故障不同，慢故障難以偵測，
因為受影響的節點仍持續回應，僅是速度較慢。
標準健康檢查（TCP 探測、心跳）
無法區分減速達 5 倍的節點與暫態負載下的健康節點。
故障訊號的模糊問題在高系統利用率下尤為明顯，
此時\emph{所有}節點都會因排隊效應而延遲升高。

\subsection{Power-of-Two-Choices 負載均衡}

Power-of-Two-Choices（P2C）~\cite{mitzenmacher2001power}
是現代分散式系統中的主流負載均衡演算法，
Envoy~\cite{envoy}、nginx 及眾多雲原生基礎設施均採用此演算法。
其運作方式為：每當請求到達時，P2C 隨機抽樣兩個候選工作節點，
並將請求路由至佇列較短的那個。

P2C 對慢故障有一種隱式緩解效果，雖然重要，一般討論卻鮮少提及。
當某節點的減速因子為 $s$ 時，其處理速度降低 $s$ 倍，
佇列隨之增長。P2C 比較佇列深度時便傾向於避開該節點，
將流量導向較健康的節點。
第~\ref{sec:nonmono}~節將展示，
此隱式緩解在極端嚴重度（$s > \sstar$）下效果顯著，
但在中度嚴重度（$s < \sstar$）下則不足，
而後者正是顯式緩解策略能發揮作用的區間。

\subsection{現有緩解方法}

\subsubsection*{對沖請求}

Dean 與 Barroso~\cite{dean2013tail}提出一種緩解尾延遲的策略：
等待一小段時間（通常為預期服務時間）後，
將同一請求的副本發送至另一台伺服器，
客戶端採用先到達的回應。
此方法能有效處理隨機性慢請求，
然而如本文後續分析所示，對持續性慢故障則無效，
因為後續請求仍會持續路由至同一故障節點。

\subsubsection*{異常值黑名單}

Cassandra 的動態偵測器（dynamic snitch）~\cite{lakshman2010cassandra}
追蹤各節點延遲，
並將延遲超出叢集中位數達 $k$ 個標準差以上的節點臨時加入黑名單。
此為前瞻式機制，能將後續流量重新導向，
但在高負載下存在振盪風險：
將節點加入黑名單實質上等於減少系統可用容量，
進而增加剩餘節點的負擔。

\subsubsection*{斷路器}

Istio 與 Envoy 內建斷路器機制，
當錯誤率或延遲超過閾值時觸發，
將流量從不健康的端點重新導向。
此類機制雖然有效，但需手動調整閾值，
且無法自適應於不同的負載條件。

\subsection{緩解陷阱}
\label{sec:trap}

推動本研究的關鍵動機之一，是\emph{緩解陷阱}（mitigation trap）：
在高系統負載下，
樸素的緩解策略\emph{反而可能適得其反}。
考慮 85\% 負載下的對沖請求：
每次對沖產生一個額外請求，
使健康節點的利用率從 $\rho = 0.85$ 推向飽和。
M/D/1 模型的排隊延遲依 $\rho/(2(1-\rho))$ 遞增，
形成正回饋迴圈：
對沖增加負載，負載增加延遲，延遲觸發更多對沖。

實驗結果驗證了此現象（第~\ref{sec:eval-lit}~節），
並據此歸納出本文負載感知自適應框架的核心設計原則：
任何緩解策略皆須將自身的容量成本納入考量。
